\documentclass[10pt,a4paper]{amsart}
\usepackage[margin=19mm]{geometry}
\usepackage{amsmath,amssymb,mathtools}
\usepackage[scr=rsfs,cal=euler]{mathalfa}
\usepackage{xfrac}
\usepackage[unicode,hidelinks,bookmarks=false]{hyperref}
\newcommand{\ie}{i.e.\ }
\newcommand{\cf}{cf.\ }
\newcommand{\resp}{respectively\ }
\newcommand{\Subsection}{\subsection}
\providecommand{\nobreakdash}{\nobreak\hbox{-}}
\setlength{\emergencystretch}{2em}
\multlinegap0pt
\usepackage{bm}
\usepackage{bbm}
\usepackage{dsfont}
\usepackage{xspace}
\usepackage{cancel}
\usepackage{mathtools}
\usepackage{amsthm}
\makeatletter
\@ifundefined{theorem}{\newtheorem{theorem}{Theorem}}{}
\@ifundefined{lemma}{\newtheorem{lemma}[theorem]{Lemma}}{}
\makeatother
\title[On the Piecewise Linear Perturbations of the Doubling Map]{On the Piecewise Linear Perturbations of the Doubling Map}
\author{Anubrato Bhattacharyya, Kuntal Banerjee}
\date{Source: arXiv:2510.11441v1 (CC BY 4.0)}
\begin{document}
\maketitle

\noindent\emph{Source and licence.} On the Piecewise Linear Perturbations of the Doubling Map. Version arXiv:2510.11441v1 (CC BY 4.0), distributed under the Creative Commons Attribution 4.0 International licence, as recorded in the arXiv licence metadata for this exact version and shown on the abstract page https://arxiv.org/abs/2510.11441v1. Full rights notice: licenses/arxiv-2510.11441-CC-BY-4.0.txt.\par\smallskip

\noindent\textit{Editorial scope.} Counted unit: the Lemma whose source label is \texttt{setting up1} in the article (source0.tex lines 629--642, source label \texttt{setting up1}), delivered together with the article's own complete original proof of it (source0.tex lines 643--681, one \texttt{proof} environment of 39 lines). No mathematics is written, repaired or re-derived: statement and proof are the article's own text line for line. Required context retained uncounted: none. Every definition, display and label of those lines is retained unchanged. Omitted from this package and not redistributed: 1--628, 682--1157. Nothing inside a retained excerpt is omitted or altered: every retained body is delivered exactly as the article's lines are written, so no omission receipt of any kind accompanies this package. Citations: the retained excerpts carry no citation command at all, so no bibliography entry of the article is transcribed and no reference is redistributed. Reference adaptation: none was needed and none was made; every reference inside the delivered unit resolves inside it, so no unresolved reference or citation is produced. Printed slips kept literal: no printed word, symbol, hypothesis or number of the counted statement or of its proof was altered. Layout and class: the article's own class is \texttt{amsart} with packages including amssymb, amsmath, amsthm, fullpage, stmaryrd, mathtools, soul, caption, color, subcaption, float, graphicx, tabulary, hyperref; the wrapper is the lane's pinned \texttt{amsart} editorial wrapper at 10pt on A4, which restates the theorem-like environments the retained text needs and adds the article's own macro definitions that the retained text needs (none), with extra wrapper packages (bm, bbm, dsfont, xspace, cancel, mathtools). The delivered numbering is the unit's own. Assets: no retained span contains a figure environment, a graphics-inclusion command, a PDF-page inclusion, a table or a TikZ picture; no image file is copied into this package, the package declares no asset dependency and no image file is redistributed with it. Licence: version arXiv:2510.11441v1 of this article is distributed under the Creative Commons Attribution 4.0 International licence, as recorded in the arXiv licence metadata for this exact version and shown on the abstract page https://arxiv.org/abs/2510.11441v1. A full rights notice is delivered at licenses/arxiv-2510.11441-CC-BY-4.0.txt. Characters: every retained excerpt is pure ASCII, so the delivered engine, XeLaTeX with its default Latin Modern text font, reports no missing character.
\begin{lemma}\label{setting up1}
    Let $b\lesssim1$ be such that there is some $(a,x)\in \text{SM}^4(b)$ with $\dfrac{2-b}{2}<a <1/2+b/2$ and $F_{a,b}(I_-)\subset (2,2.5)$.
    Pick such an $(a,x)\in \text{SM}^{4}(b)$ and fix it. 
    \begin{enumerate}
        \item Let $\alpha \in \mathbb{T}$ so that $|\alpha-a| \leq \text{ } f_{a,b}(x) \wedge (1/2-f_{a,b}(x))$ then $f_{\alpha,b}(x)\in I_+$.
        \vspace{.3cm}
        
        \item Let $\alpha \in \mathbb{T}$ so that $\alpha$ satisfies the bound in (1) (i.e. $F_{a,b}(x) \in (2,2.5) $ ) 
        and $|\alpha-a| \leq \text{ } \dfrac{f^2_{a,b}(x) \wedge (1/2-f^2_{a,b}(x))}{B^++1} $ then $f^2_{\alpha,b}(x)\in I_+$.
          \vspace{.3cm}
        \item Let $\alpha \in \mathbb{T}$ so that $\alpha$ satisfies bounds in  (1), (2) and  $|\alpha-a| \leq \text{ } \dfrac{f^3_{a,b}(x) \wedge (1/2-f^3_{a,b}(x))}{B^{+^{2}}+B^++1} $ then $f^3_{\alpha,b}(x)\in I_+$.
        
    \end{enumerate}
\end{lemma}

\begin{proof}
Observe that $0<F_{a,b}(x)<2.5$ for $x\in I$ and $0<F_{a,b}(A^+)< 2.5$ under the assumed conditions. For $\alpha \in \mathbb{T}$, denote $\alpha^+=\alpha-b/2, \alpha^-=\alpha+3b/2$, then the following hold. Since $\dfrac{2-b}{2}\geq b/2$ if $b \in [0,1] $, we get that $b/2<a<1/2+b/2$. 
 
 \vspace{.3cm}
 
\begin{enumerate}
    \item Under the assumptions on $b$, $f_{a,b}(I_-) \subset I_+$ and  $F_{a,b}(x) \in (2,2.5)$ . Then  $F_{a,b}(x)-2=f_{a,b}(x)$ and $2.5-F_{a,b}(x)=0.5-f_{a,b}(x)$.
    Since $F_{a,b}(x)-F_{\alpha,b}(x)=a-\alpha$ and (1) follows.
    
    \vspace{.3cm}
    \item Note that $f^2_{a,b}(x) \in I_+$ if,  $F^2_{a,b}(x)  \in (4,4.5 ) \bigcup (5,5.5)\bigcup (6,6.5)$, $|F_{a,b}^2(2,2.5)|= 2-(1-b)$.
    \vspace{.1cm}
\begin{itemize}
    \item Now $F^2_{a,b}(x)  \in (4,4.5 )$ and if  $|\alpha-a|\leq \dfrac{(4.5 -F^2_{a,b}(x) ) \wedge (F^2_{a,b}(x)-4)  }{(B^++1)}$ then,\\ $F^2_{\alpha,b}(x) \in (4,4.5) $. Observe that,
    \vspace{.3cm}
    \begin{itemize}
        \item[]  $F^2_{a,b}(x)  \in (4,4.5 ) \implies 4.5 -F^2_{a,b}(x) = 1/2-f_{a,b}^2(x) \text{ as } F^2_{a,b}(x)-4 = f^2_{a,b}(x) $.
        \vspace{.3cm}
        \item[] So, $\dfrac{(4.5 -F^2_{a,b}(x))  \wedge( F^2_{a,b}(x)-4)  }{(B^++1)}=\dfrac{(1/2 -f^2_{a,b}(x))  \wedge f^2_{a,b}(x) }{(B^++1)}.$ 
    \end{itemize}
   \vspace{.3cm}
    \item  We can similarly deal with $F^2_{a,b}(x)  \in (5,5.5 ) $ and $F^2_{a,b}(x)  \in (6,6.5 ) $
    
\end{itemize}
    Now say $F_{a,b}^2(x) \in (4,4.5)$ and let $\alpha \in \mathbb{T} $ such that it satisfies (1) and
    \vspace{.3cm}
    
    $|\alpha-a| \leq \text{ } \dfrac{f^2_{a,b}(x) \wedge (1/2-f^2_{a,b}(x))}{B^++1}$. Because $\alpha$ satisfies (1), we have $F_{\alpha,b}(x) \in (2,2.5)$. Hence, 
    $F^2_{a,b}(x)=F_{a,b}(B^-x+A^-)=B^+(B^-x+A^--2) +A^++4=$ and $F^2_{\alpha,b}(x)=F_{\alpha,b}(B^-x+\alpha^-)=B^+(B^-x+\alpha^--2) +\alpha^++4$, or, $F^2_{\alpha,b}-F^2_{a,b}= (B^++1)(\alpha-a)$. We can Deal with the other cases similarly.
    
    \vspace{.3cm}

    \item Now $f_{a,b}^3(x) \in I_+$ means, \[
    F_{a,b}^3(x) \in \bigcup_{j=8}^{14}(j,j+1/2).
    \]
    Without loss of generality consider, $F_{a,b}^3(x) \in (8,8.5)$. Assume $\alpha \in \mathbb{T}$ satisfies (1),(2) then as before, $F_{a,b}^3(x)-F_{\alpha,b}^3(x)=(B^{+{2}}+B^++1)(a-\alpha)$. Similar considerations for other cases makes the proof (3) complete.
\end{enumerate}
 
\end{proof}

\end{document}
